% P15-U1: maximal inequalities and the complete three-series criterion.
\chapter{Independent Random Series}
\label{sec:p15-independent-random-series}

The convergence of an infinite random sum is a statement about all its
late partial sums. A maximal inequality controls those sums together;
truncation then separates large jumps, deterministic drift and random
fluctuation. The three-series theorem gives an exact criterion for
independent real summands, without assuming their original expectations
or variances exist.

The source is Durrett's author Version 5 dated 11 January 2019,
Theorems 2.5.5--2.5.6, Example 2.5.7 and Theorem 2.5.8
\cite[pp.~84--85]{DurrettPTEAuthor2019}. The necessity proof below is
elementary and complete. The source's different necessity argument,
deferred to Example 3.4.12, is also given in the annex with its exact
Lindeberg--Feller premise displayed separately
\cite[pp.~148--151]{DurrettPTEAuthor2019}.

\section{Partial sums and simultaneous bounds}
\label{sec:p15-series-partial-sums}

\begin{hbdefinition}{def:p15-random-series}{Partial sums, finite convergence and tail oscillation}
Let \(X_1,X_2,\ldots\) be real measurable random variables on
\((\Omega,\mathcal F,\mathbb P)\). Set
\[
 S_0=0,\qquad S_n=\sum_{j=1}^nX_j,\qquad
 C_X=\{\omega:(S_n(\omega))_{n\geq0}
                         \text{ converges in }\mathbb R\}.
\]
Thus convergence of \(\sum_nX_n\) always means a finite real limit.
Almost-sure convergence means \(\mathbb P(C_X)=1\).
Absolute convergence means \(\sum_n|X_n|<\infty\).
Conditional convergence means finite convergence with failure of absolute
convergence. A real random variable is centered when its expectation
exists and equals zero.
For each integer \(M\geq0\), the extended nonnegative random variable
\[
 w_M=\sup_{m,n\geq M}|S_m-S_n|
\]
is the tail oscillation. Independence of a sequence means independence
of every finite subfamily, as in
Definition~\ref{def:p15-independent-random-elements}.
\index{random series}\index{partial sum}\index{tail oscillation}
\end{hbdefinition}

\begin{hbtheorem}{thm:p15-series-event-calculus}{proved=annex:p15.series-event-calculus}{The measurable Cauchy event}
The convergence event has the countable description
\[
 C_X=\bigcap_{r\geq1}\ \bigcup_{M\geq0}
       \bigcap_{m,n\geq M}\{|S_m-S_n|\leq1/r\}
     =\{w_M\longrightarrow0\}.
\]
It belongs to the tail sigma-algebra
\(\bigcap_{k\geq1}\sigma(X_k,X_{k+1},\ldots)\).
On this event \(X_n\to0\). Changing finitely many summands at an
outcome preserves convergence and absolute convergence there.
If the summands are independent, then
\(\mathbb P(C_X)\in\{0,1\}\).
\end{hbtheorem}

\begin{hbtheorem}{thm:p15-kolmogorov-maximal}{proved=annex:p15.kolmogorov-maximal}{Kolmogorov's maximal inequality}
Let \(n\geq1\) and let \(X_1,\ldots,X_n\) be independent, with
\(\mathbb EX_j=0\) and \(\mathbb EX_j^2<\infty\).
For every real \(x>0\),
\[
 \mathbb P\!\left(\max_{1\leq k\leq n}|S_k|\geq x\right)
 \leq \frac{\operatorname{Var}(S_n)}{x^2}
 =\frac{1}{x^2}\sum_{j=1}^n\operatorname{Var}(X_j).
\]
The maximum is over the whole initial block; a bound for \(S_n\)
alone would not imply this assertion.
\end{hbtheorem}

\begin{hbtheorem}{thm:p15-variance-series-convergence}{proved=annex:p15.variance-series-convergence}{Summable variances give convergence}
Let \((X_n)_{n\geq1}\) be independent, centered real random variables
with finite variances. If
\[
 V=\sum_{n\geq1}\operatorname{Var}(X_n)<\infty,
\]
then \(\sum_nX_n\) converges almost surely to a finite random
variable \(S_\infty\). Define its value to be zero off the
convergence event. This version has finite second moment and
\[
 \mathbb ES_\infty=0,\qquad
 \mathbb E(S_\infty-S_M)^2
      =\sum_{j>M}\operatorname{Var}(X_j),\qquad
 \operatorname{Var}(S_\infty)=V.
\]
In particular \(S_M\to S_\infty\) in mean square, meaning that the
displayed mean squared difference tends to zero.
\end{hbtheorem}

\section{The bounded converse and three-series theorem}
\label{sec:p15-three-series}

\begin{hblemma}{thm:p15-bounded-fourth-moment}{proved=annex:p15.bounded-fourth-moment}{A fluctuation bound for bounded centered sums}
Let \(U_1,\ldots,U_n\) be independent and centered, with
\(|U_j|\leq M\) almost surely for a fixed finite \(M\geq0\).
Set \(W=\sum_jU_j\) and \(v=\sum_j\mathbb EU_j^2\).
Then
\[
 \mathbb EW^2=v,\qquad
 \mathbb EW^4\leq M^2v+3v^2.
\]
If \(v>0\), then
\[
 \mathbb P\!\left(|W|>\sqrt{v/2}\right)
       \geq\frac{v}{4(M^2+3v)}.
\]
Consequently, for an independent centered sequence uniformly bounded
by \(M\), almost-sure convergence of its series implies summability
of its variances.
\end{hblemma}

\begin{hblemma}{thm:p15-series-independent-copy}{proved=annex:p15.series-independent-copy}{An independent copy and symmetrization}
Let \((Y_n)\) be independent and satisfy \(|Y_n|\leq A\) almost
surely for some finite \(A\geq0\). On the finite product probability
space \((\Omega^2,\mathcal F\otimes\mathcal F,
\mathbb P\otimes\mathbb P)\), define
\[
 \widehat Y_n(\omega,\omega')=Y_n(\omega),\qquad
 \widetilde Y_n(\omega,\omega')=Y_n(\omega'),\qquad
 Z_n=\widehat Y_n-\widetilde Y_n.
\]
The two copied sequences together form an independent family.
The variables \(Z_n\) are independent, centered and symmetric in
law, where symmetry means \(Z_n\) and \(-Z_n\) have the same law.
Moreover,
\[
 |Z_n|\leq2A\quad\hbox{a.s.},\qquad
 \operatorname{Var}(Z_n)=2\operatorname{Var}(Y_n).
\]
If \(\sum_nY_n\) converges almost surely, so does \(\sum_nZ_n\).
\end{hblemma}

\begin{hbtheorem}{thm:p15-bounded-series-criterion}{proved=annex:p15.bounded-series-criterion}{The exact criterion for uniformly bounded summands}
For independent real \(Y_n\) with a common finite constant \(A\geq0\)
and almost-sure bounds \(|Y_n|\leq A\), the series \(\sum_nY_n\) converges almost surely
if and only if both
\[
 \sum_{n\geq1}\mathbb EY_n\quad\text{converges in }\mathbb R,
 \qquad
 \sum_{n\geq1}\operatorname{Var}(Y_n)<\infty.
\]
The first condition permits conditional convergence of the numerical
series of means.
\end{hbtheorem}

\begin{hbdefinition}{def:p15-three-series-truncation}{Deletion truncation and the three series}
For a fixed real cutoff \(A>0\), put
\[
 Y_n^{(A)}=X_n\mathbf1_{\{|X_n|\leq A\}},\quad
 p_n^{(A)}=\mathbb P(|X_n|>A),\quad
 \mu_n^{(A)}=\mathbb EY_n^{(A)},\quad
 v_n^{(A)}=\operatorname{Var}(Y_n^{(A)}).
\]
A value outside \([-A,A]\) is replaced by zero; values at either
endpoint are retained. This is deletion truncation, rather than clipping
to an endpoint. The truncated variables have finite means and variances
even when the original variables do not. The three series are
\(\sum_np_n^{(A)}\), \(\sum_n\mu_n^{(A)}\) and
\(\sum_nv_n^{(A)}\).
\index{three-series theorem}\index{truncation!random series}
\end{hbdefinition}

\begin{hbtheorem}{thm:p15-three-series}{proved=annex:p15.three-series}{Kolmogorov's three-series theorem}
Let \((X_n)\) be independent real random variables and fix any
\(A>0\). Then \(\sum_nX_n\) converges almost surely if and only if
all three conditions hold:
\[
 \sum_np_n^{(A)}<\infty,\qquad
 \sum_n\mu_n^{(A)}\text{ converges in }\mathbb R,\qquad
 \sum_nv_n^{(A)}<\infty.
\]
No integrability hypothesis is imposed on the untruncated \(X_n\).
The criterion holds at one positive cutoff if and only if it holds at
every positive cutoff. If any condition fails at a fixed positive
cutoff, the event of finite convergence has probability zero.
\end{hbtheorem}

\section{A separately stated central-limit route}
\label{sec:p15-series-lindeberg-route}

\begin{hbdefinition}{def:p15-series-lindeberg-premise}{The Lindeberg--Feller assertion used by the source}
Let \(\Phi(t)=\mathbb P(G\leq t)\), where \(G\) has the standard
normal law \(N(0,1)\) of Definition~\ref{def:p15-normal-law}.
The following assertion is denoted \(\mathrm{LF}\):
for every triangular array of real random variables
\(U_{n,m}\), \(1\leq m\leq n\), independent within each row,
with \(\mathbb EU_{n,m}=0\) and finite second moments, if for a
finite real \(\sigma>0\)
\[
 \sum_{m=1}^n\mathbb EU_{n,m}^2\longrightarrow\sigma^2>0,
 \qquad
 \sum_{m=1}^n
 \mathbb E\!\left[U_{n,m}^2\mathbf1_{\{|U_{n,m}|>\varepsilon\}}\right]
       \longrightarrow0\quad\text{for every }\varepsilon>0,
\]
then, for every \(t\in\mathbb R\),
\[
 \mathbb P\!\left(\sum_{m=1}^nU_{n,m}\leq t\right)
                   \longrightarrow\Phi(t/\sigma).
\]
Independence between different rows is not required. This is the
distribution-function form of Durrett's Theorem 3.4.10
\cite[pp.~148--149]{DurrettPTEAuthor2019}; its proof is reserved for
the central-limit chapter in the source. The complete proof is supplied
in \cref{thm:p15-clt-lindeberg-feller}, with its distinct moment route
in \cref{thm:p15-clt-moment-array-proof}.
It is the premise used by the next source argument, and is not a premise of
Theorem~\ref{thm:p15-three-series}.
\end{hbdefinition}

\begin{hbtheorem}{thm:p15-three-series-source-route}{proved=annex:p15.three-series-source-route}{The source's alternative necessity argument}
Use \(\mathrm{LF}\), as stated in
Definition~\ref{def:p15-series-lindeberg-premise} and proved in
\cref{thm:p15-clt-lindeberg-feller}.
For independent \(X_n\) whose series converges almost surely, the
three conditions of Theorem~\ref{thm:p15-three-series} also follow
by the following central-limit argument. After proving the large-jump
condition, put \(Y_n=Y_n^{(A)}\) and
\(c_n=\sum_{m=1}^n\operatorname{Var}(Y_m)\).
If \(c_n\to\infty\), the centered rows
\((Y_m-\mathbb EY_m)/\sqrt{c_n}\) meet \(\mathrm{LF}\), while
\(\sum_{m\leq n}Y_m/\sqrt{c_n}\to0\) almost surely.
These conclusions would force the deterministic numbers
\(-\sum_{m\leq n}\mathbb EY_m/\sqrt{c_n}\) to have distribution
functions tending to \(\Phi\), a contradiction. Thus the variances
are summable, and convergence of the mean series follows by subtraction.
The annex supplies all steps, including the distribution-function
perturbation argument and the treatment of any initial zero \(c_n\).
\end{hbtheorem}

\section{Calculations and hypothesis boundaries}
\label{sec:p15-series-examples}

\begin{hbtheorem}{thm:p15-signed-power-series}{proved=annex:p15.signed-power-series}{The exact signed-power threshold}
Let \((\varepsilon_n)\) be independent \(\{-1,1\}\)-valued random variables with
\(\mathbb P(\varepsilon_n=1)=\mathbb P(\varepsilon_n=-1)=1/2\),
and let \(\alpha\in\mathbb R\). Then
\[
 \sum_{n\geq1}\varepsilon_n n^{-\alpha}
       \text{ converges a.s.}\quad\Longleftrightarrow\quad\alpha>1/2.
\]
It converges absolutely exactly when \(\alpha>1\), and hence
converges conditionally almost surely for \(1/2<\alpha\leq1\).
For \(\alpha\leq1/2\) its finite-convergence event has probability
zero. This completes the classification of the source's Example 2.5.7.
\end{hbtheorem}

\begin{hbtheorem}{thm:p15-series-boundary-examples}{proved=annex:p15.series-boundary-examples}{Separating drift, large jumps, variance and independence}
The following examples use cutoff \(A=1\).
\begin{enumerate}
\item For deterministic \(X_n=(-1)^{n+1}/n\), the large-jump and
variance series are zero and the mean series converges conditionally.
Thus absolute summability of the means is unnecessary.
\item Each of the three conditions can fail while the other two hold.
For independent \(X_n\) with masses \(1/(2n)\) at each of
\(n+1\) and \(-(n+1)\), and mass \(1-1/n\) at zero, only the
large-jump series fails. For deterministic \(X_n=1/n\), only the
mean series fails. For \(X_n=\varepsilon_n/\sqrt n\), only the
variance series fails. All three random series diverge almost surely.
\item Independent variables with masses \(2^{-n-1}\) at each of
\(2^n\) and \(-2^n\), and the remaining mass at zero, give an
almost-surely absolutely convergent series although
\(\sum_n\operatorname{Var}(X_n)=\sum_n2^n=\infty\).
The bounded converse cannot be extended to unbounded summands this way.
\item With a single fair sign \(\varepsilon\), the dependent
centered variables \(X_n=\varepsilon/n\) have summable variances
but their series diverges at every outcome. For the maximal inequality,
the dependent centered pair \(X_1=\varepsilon,X_2=-\varepsilon\)
has terminal variance zero and maximum absolute partial sum one.
The uncentered one-term variable \(X_1=1\) also has variance zero
and absolute partial sum one. Thus independence and centering are
substantive hypotheses of the stated inequalities and convergence result.
\end{enumerate}
All independent sequences with the specified discrete laws exist by
Theorem~\ref{thm:p15-independent-sequence-existence}.
\end{hbtheorem}

\begin{proofboundary}{def:p15-series-next-boundary}{The next probability results}
This chapter proves the maximal inequality, convergence from summable
variances, and both directions of the three-series theorem. It also
gives the full source necessity route of Example 3.4.12. Its
Lindeberg--Feller premise is supplied in
\cref{thm:p15-clt-lindeberg-feller}, and the resulting unconditional
source construction is also \cref{thm:p15-clt-three-series-source}.
Kronecker's lemma and the strong law derived
from it are supplied in
\cref{thm:p15-kronecker-lemma,thm:p15-iid-series-strong-law},
retaining the separate source proof beginning with Theorem 2.5.9
\cite[pp.~85--87]{DurrettPTEAuthor2019}.
\end{proofboundary}
